%% kapa: exam=1 round=35 year=2024 period=1 form=A subject=real_estate subject_ko=부동산학원론 seq=28 number=108
%% source: Q-Net 감정평가사 기출문제 (cst00309) artlSeq=5213229 file=raw/1008/5213229/1교시 문제지 원본.pdf
%% conversion: kordoc 4.24.51 -> tools/gen1.py (auto); review=proofread-sonnet; flags=none
\begin{question}{108}
부동산투자회사법령상 부동산투자회사에 관한 내용으로 옳지 않은 것은?
\begin{choices}
  \item 영업인가를 받거나 등록을 한 날부터 최저자본금준비기간이 지난 자기관리 부동산투자회사의 최저자본금은 70억원 이상이 되어야 한다.
  \item 최저자본금준비기간이 끝난 후에는 매 분기 말 현재 총자산의 100분의 80 이상을 부동산, 부동산 관련 증권 및 현금으로 구성하여야 한다. 이 경우 총자산의 100분의 70 이상은 부동산(건축 중인 건축물을 포함한다)이어야 한다.
  \item 부동산투자회사는 부동산 등 자산의 운용에 관하여 회계처리를 할 때에는 금융감독원이 정하는 회계처리기준에 따라야 한다.
  \item 부동산투자회사의 상근 임원은 다른 회사의 상근 임직원이 되거나 다른 사업을 하여서는 아니 된다.
  \item 위탁관리 부동산투자회사란 자산의 투자ㆍ운용을 자산관리회사에 위탁하는 부동산투자회사를 말한다.
\end{choices}
\end{question}
\answer{3}
